% Native TikZ geometry, reconstructed from the supplied reference image.
% Ambient model: p(t)=(t,a*t^2,b*t^2), x=p(0).
% T_x M=span(e_1), nu_0=e_2, Ran(P_0)=span(e_2), Ran(S_0)=span(e_3).
% Oblique display: (v1,v2,v3) -> (v1-0.25*v2+0.60*v3,v2+0.65*v3).
% The three colored arrows are illustrative WEIGHTED leading components;
% their scalar coefficients are suppressed in the graphic and stated in
% the caption. The closing remainder is deliberately nonzero.
\begingroup
\definecolor{truecurve}{HTML}{292E34}
\definecolor{yycurve}{HTML}{315F91}
\definecolor{pcomponent}{HTML}{2868A8}
\definecolor{scomponent}{HTML}{148A85}
\definecolor{bcomponent}{HTML}{C4772B}
\definecolor{fullnormal}{HTML}{525E69}
\definecolor{planeedge}{HTML}{CDD0D3}
\definecolor{guide}{HTML}{9098A1}
\def\curva{0.08}
\def\curvb{0.04}
\def\leadP{2.05}
\def\leadS{0.62}
\def\leadB{0.49}
\def\restP{0.05}
\def\restS{0.01}
% Small parallel translation of the free direction marker within the plane.
% The true P_0 component itself still starts at x and lies in Ran(P_0).
\def\pguideS{-0.20}
\pgfmathsetmacro{\offsetP}{\leadP+\restP}
\pgfmathsetmacro{\offsetS}{\leadS+\leadB+\restS}
\pgfmathdeclarefunction{tx}{1}{\pgfmathparse{#1+(-0.25*\curva+0.60*\curvb)*#1*#1}}
\pgfmathdeclarefunction{ty}{1}{\pgfmathparse{(\curva+0.65*\curvb)*#1*#1}}
\pgfmathdeclarefunction{normalfactor}{1}{\pgfmathparse{
  (2*\curva*\offsetP+2*\curvb*\offsetS)*#1 /
  (1+4*(\curva*\curva+\curvb*\curvb)*#1*#1)}}
\pgfmathdeclarefunction{qx}{1}{\pgfmathparse{
  tx(#1)-0.25*\offsetP+0.60*\offsetS
  -normalfactor(#1)*(1+(-0.50*\curva+1.20*\curvb)*#1)}}
\pgfmathdeclarefunction{qy}{1}{\pgfmathparse{
  ty(#1)+\offsetP+0.65*\offsetS
  -normalfactor(#1)*(2*\curva+1.30*\curvb)*#1}}
% Normal-plane points are constructed from their actual basis coefficients.
\newcommand{\NormalCoordinate}[3]{%
  \coordinate (#1) at ({-0.25*(#2)+0.60*(#3)},{(#2)+0.65*(#3)});}
\begin{tikzpicture}[
  x=1.40cm,y=1.40cm,line cap=round,line join=round,
  >=Latex,font=\normalsize,
  every node/.style={inner sep=1.2pt,outer sep=0pt},
  field/.style={-{Latex[length=2.2mm,width=1.5mm]},line width=1.25pt},
  leader/.style={draw=guide,line width=0.38pt},
  axisguide/.style={draw=guide,line width=0.55pt,-{Latex[length=1.4mm,width=0.9mm]}}
]
\path[use as bounding box] (-3.95,-1.02) rectangle (4.28,4.56);
\NormalCoordinate{X}{0}{0}
\NormalCoordinate{A}{\leadP}{0}
\NormalCoordinate{B}{\leadP}{\leadS}
\NormalCoordinate{C}{\leadP}{\leadS+\leadB}
\NormalCoordinate{Q}{\offsetP}{\offsetS}
\NormalCoordinate{Fll}{-0.45}{-0.65}
\NormalCoordinate{Ful}{2.60}{-0.65}
\NormalCoordinate{Fur}{2.60}{1.35}
\NormalCoordinate{Flr}{-0.45}{1.35}
\NormalCoordinate{PaxisLow}{-0.33}{\pguideS}
\NormalCoordinate{PaxisHigh}{2.62}{\pguideS}
\NormalCoordinate{SaxisLow}{0}{-0.50}
\NormalCoordinate{SaxisHigh}{0}{1.22}

% The affine true normal plane: very light gray fill, thin gray outline.
\filldraw[fill=black!1.5,draw=planeedge,line width=0.35pt]
  (Fll)--(Ful)--(Fur)--(Flr)--cycle;

% Tangent and the two orthogonal ambient normal directions (oblique view).
\draw[guide,line width=0.48pt,dash pattern=on 2.6pt off 2.1pt,
  -{Latex[length=1.5mm,width=1mm]}] (-3.50,0)--(3.48,0);
\draw[axisguide] (PaxisLow)--(PaxisHigh);
\draw[axisguide] (SaxisLow)--(SaxisHigh);

% Curves are analytic, not unrelated hand-drawn Bezier paths.
\draw[truecurve,line width=1.35pt]
  plot[domain=-3.40:3.40,samples=161,variable=\t] ({tx(\t)},{ty(\t)});
% q(t)=p(t)+Q(t)U, where Q(t) is the true normal projector at p(t).
% Thus the blue curve is a genuine normal graph and q(0)=x+u_{r_1}(x).
\draw[yycurve,line width=1.35pt]
  plot[domain=-4.60:4.60,samples=201,variable=\t] ({qx(\t)},{qy(\t)});

% Independent full normal displacement, with its own dashed encoding.
\draw[fullnormal,line width=0.85pt,dash pattern=on 3.3pt off 2.2pt,
  -{Latex[length=2.1mm,width=1.3mm]}]
  (X)--(Q);

% Head-to-tail leading components. B-A and C-B are exactly collinear.
\draw[field,pcomponent] (X)--(A);
\draw[field,scomponent] (A)--(B);
\draw[field,bcomponent] (B)--(C);
% C is close to, but not exactly Q. The remaining O(sigma^4) term is
% suppressed at this drawing scale and is explained in the caption.

% Point markers indicate positions; the colored paths indicate vectors.
\fill[truecurve] (X) circle[radius=1.8pt];
\fill[yycurve] (Q) circle[radius=1.9pt];
\node[below=4pt,text=truecurve] at (X) {$x$};

% Direct component labels, with thin neutral leaders where needed.
\node[anchor=east,text=pcomponent] at (-0.58,1.08) {$P_0H$};
\draw[leader] (-0.56,1.08)--(-0.25,1.00);
\node[anchor=west,text=scomponent] at (0.82,2.05) {$S_0H$};
\draw[leader] (0.79,2.05)--($(A)!0.50!(B)$);
\node[anchor=west,text=bcomponent] at (0.84,2.56) {$b_\perp$};
\draw[leader] (0.80,2.56)--($(B)!0.50!(C)$);

% nu_0 identifies the P_0 axis, not an additional bias displacement.
\node[anchor=east,text=fullnormal] at (-0.92,2.58)
  {$\nu_0\; (\operatorname{Ran}P_0)$};
\node[anchor=west,text=fullnormal] at (0.91,0.83) {$\operatorname{Ran}S_0$};
\node[anchor=west,text=guide] at (3.58,-0.02) {$T_x\mathcal M$};
\node[anchor=west,text=fullnormal] at (0.94,1.37) {$u_{r_1}(x)$};
\draw[leader] (0.90,1.37)--($(X)!0.485!(Q)$);

% Position and manifold labels are separated from vector labels.
\node[anchor=south,text=yycurve] at (-0.40,3.62) {$x+u_{r_1}(x)$};
\draw[leader] (-0.40,3.58)--(Q);
\node[anchor=west,text=fullnormal] at (1.02,4.13) {$x+N_x\mathcal M$};
\draw[leader] (0.96,4.12)--(Fur);
\node[anchor=west,text=yycurve] at (3.60,3.85) {$\mathcal M_{r_1}^{\mathrm{YY}}$};
\draw[leader] (3.57,3.85)--({qx(4.60)},{qy(4.60)});
\node[anchor=west,text=truecurve] at (3.62,1.17) {$\mathcal M$};
\draw[leader] (3.59,1.17)--({tx(3.40)},{ty(3.40)});
\end{tikzpicture}
\endgroup
